\expandafter\def\csname HistoricalNoteText391\endcsname{Ein Mangel in der festen Verbindung der Theile des schwingenden Apparats unter einander hat immer eine schnelle unregelmässige Abnahme der Schwingungsbögen zur Folge.}
\expandafter\def\csname HistoricalNoteText382\endcsname{Dass zu einer Zeit, wo die Declination schnell wechselnde starke Änderungen erleidet, \textit{sehr kleine} Schwingungen (die aber schon an sich zur Bestimmung der Schwingungsdauer wenig tauglich sind) eine ganz entstellte Dauer zeigen können, braucht kaum bemerkt zu werden.}
\SourceLeaf{382}allerdings die Schwingungsdauer von veränderlichen Elementen abhängig, und daher auch, selbst streng genommen, beständigen Veränderungen unterworfen: allein von ganz ungewöhnlichen Fällen abgesehen\HistoricalNote{382}, wird diese Veränderlichkeit während einer nicht ganz kleinen Zeit als ganz unmerklich zu betrachten sein, so wie jedenfalls der aus einer beträchtlichen Anzahl von Schwingungen abgeleitete Werth der Dauer Einer Schwingung, dem Mittelwerthe der einzelnen in Betracht kommenden Elemente während dieser Zeit entsprechen wird. Es ist aber offenbar gar nicht nöthig, den Bewegungen der Nadel während eines solchen Zeitraumes ununterbrochen zu folgen, sondern es reicht hin, die Zeiten der ersten und der letzten Elongation zu kennen, so bald man von der Schwingungsdauer einen so weit genäherten Werth besitzt, dass über die \textit{Anzahl} der Schwingungen während des verflossenen Zeitraums kein Zweifel übrig bleiben kann, was dadurch noch erleichtert wird, dass man allemal (nach der Gleichnamigkeit oder Ungleichnamigkeit der ersten und letzten Elongation) im Voraus weiss, ob diese Anzahl gerade oder ungerade ist.

Wenn die Schwingungsdauer nicht gar zu klein ist, so können zwischen den Vorübergängen auch die Elongationen selbst (nemlich die äussersten Scalentheile) mit aufgezeichnet werden, um daraus die zur schärfern Berechnung der Schwingungsdauer nöthigen Amplituden ableiten zu können, deren successive Abnahme überdies an sich zu merkwürdigen Betrachtungen Anlass gibt.

Die Behandlung der Beobachtungen selbst, um Resultate aus ihnen zu gewinnen, wird sich am besten an einem Beispiele zeigen lassen, wozu hier Beobachtungen am Magnetometer der Sternwarte mit fünfundzwanzigpfündigem Stabe, vom 29. November 1835 gewählt werden. Die folgende Tafel I. enthält zuerst die rohen Beobachtungen.

\begin{center}
\begingroup\fontsize{9}{11}\selectfont\setlength{\tabcolsep}{2pt}
\begin{tabular}{r|r||r|r||r|r||r|r}
& $1755.1$ && $497.8$ && $645.9$ && $1232.1$ \\[-4pt]
$21^{\mathrm{h}}55'26''9$ && $23^{\mathrm{h}}38'49''2$ && $1^{\mathrm{h}}10'12''6$ && $2^{\mathrm{h}}49'19''7$ & \\[-4pt]
& $266.0$ && $1502.2$ && $1341.5$ && $775.9$ \\[-4pt]
$56\quad8.4$ && $39\quad31.5$ && $54.2$ && $50\quad1.5$ & \\[-4pt]
& $1751.8$ && $500.1$ && $647.3$ && $1231.0$ \\[-4pt]
$51.2$ && $40\quad13.6$ && $11\quad37.0$ && $44.1$ & \\[-4pt]
& $268.5$ && $1499.1$ && $1339.4$ && $776.4$ \\[-4pt]
$57\quad33.0$ && $56.0$ && $12\quad18.4$ && $51\quad25.8$ & \\[-4pt]
& $1748.9$ && $502.6$ && $648.7$ && $1228.7$ \\[-4pt]
$58\quad15.5$ && $41\quad38.1$ && $13\quad1.3$ && $52\quad8.5$ & \\[-4pt]
& $271.6$ && $1496.5$ && $1337.0$ && $778.0$ \\[-4pt]
$57.4$ && $42\quad20.3$ && $43.0$ && $50.0$ & \\[-4pt]
& $1744.2$ && $506.0$ && $650.7$ && $1227.0$
\end{tabular}\endgroup
\end{center}

\SourceLeaf{383}Diese Beobachtungen bestehen, wie man sieht, aus vier Sätzen; die erste Columne enthält die Zeiten der Vorübergänge des Scalenpunktes 1000, die zweite die Elongationspunkte, diesmal mit der Elongation anfangend, die dem ersten Vorübergange voranging, und mit derjenigen schliessend, die auf den letzten Vorübergang folgte. Wenn man die Elongationen nicht mit beobachtet, so thut man wohl, bei jedem Vorübergange anzumerken, ob wachsende oder abnehmende Zahlen durchgingen; nach der in Göttingen befolgten Art so:
\begin{center}
\begin{tabular}{l}
$21^{\mathrm{h}}55'26''9-$ \\
$\phantom{21^{\mathrm{h}}}56\quad8.4+$ \\
u. s. f. \\
$23^{\mathrm{h}}38'49''2+$ \\
u. s. w.
\end{tabular}
\end{center}

Dies ist deswegen nöthig, um unterscheiden zu können, welche der aus den Vorübergängen abgeleiteten Elongationszeiten sich auf Minima oder Maxima beziehen. Bei der Zählung der Elongationszeiten haben wir die Gewohnheit angenommen, die erstern durch gerade, die andern durch ungerade Zahlen zu bezeichnen. Es wird daher der aus den beiden ersten Vorübergängen abgeleiteten Elongationszeit $21^{\mathrm{h}}55'47''65$ die Zahl 0 vorgesetzt u. s. f.

Die folgende Tafel II. enthält nun die nächsten aus den unmittelbaren Beobachtungen berechneten Resultate.
\begin{center}
\begingroup\fontsize{9}{11}\selectfont\setlength{\tabcolsep}{4pt}
\begin{tabular}{r|r|r||r|r|r}
0 & $21^{\mathrm{h}}55'47''65$ & $21^{\mathrm{h}}55'47''65$ & 277 & $1^{\mathrm{h}}10'33''40$ & $21^{\mathrm{h}}55'49''54$ \\
1 & $56\quad29.80$ & $47.62$ & 278 & $11\quad15.60$ & $49.56$ \\
2 & $57\quad12.10$ & $47.74$ & 279 & $11\quad57.70$ & $49.48$ \\
3 & $57\quad54.25$ & $47.71$ & 280 & $12\quad39.85$ & $49.45$ \\
4 & $58\quad36.45$ & $47.73$ & 281 & $13\quad22.15$ & $49.57$ \\[7pt]
147 & $23\quad39\quad10.35$ & $21\quad55\quad49.89$ & 418 & $2\quad49\quad40.60$ & $21\quad55\quad49.36$ \\
148 & $39\quad52.55$ & $49.91$ & 419 & $50\quad22.80$ & $49.38$ \\
149 & $40\quad34.80$ & $49.98$ & 420 & $51\quad4.95$ & $49.35$ \\
150 & $41\quad17.05$ & $50.05$ & 421 & $51\quad47.15$ & $49.37$ \\
151 & $41\quad59.20$ & $50.02$ & 422 & $52\quad29.25$ & $49.29$
\end{tabular}\endgroup
\end{center}

In der zweiten Columne stehen hier die sich ergebenden Elongationszeiten. Die Bezifferung, in der ersten Columne, hat man, für die fünf ersten von selbst; für die spätern findet sie sich auf folgende Art.


\SourceLeaf{384}Die Vergleichung der. Elongation 0 mit 4 gibt als genäherten Werth der Schwingungsdauer $42''20$; dividirt man damit die Zwischenzeit zwischen der Elongation 4 und der nächstfolgenden, $1^{\mathrm{h}}40'33''90$, und erinnert sich, dass die Ordnungszahl der letztern eine ungerade sein muss, so lässt der Quotient $142.983$ keinen Zweifel übrig, dass zwischen jenen beiden Elongationen 143 Schwingungen verflossen sein müssen; denn in der That, wollte man 141 oder 145 annehmen, so würde die Schwingungsdauer $42''7936$ oder $41''6131$ sich ergeben, viel zu stark von dem genäherten Werthe $42''20$ abweichend, um zulässig zu sein. Von 143 Schwingungen ausgehend, findet man die Schwingungsdauer $42''1951$, die man bei dem Übergange zu dem folgenden Beobachtungssätzen zum Grunde legen könnte, um ihre Bezifferung zu erhalten, obwohl in dem gegenwärtigen Falle, wo keine \textit{sehr} langen Unterbrechungen vorkommen, auch schon der erste genäherte Werth überall ausreicht.

Um die Schwingungsdauer genauer zu erhalten, und selbst ihre Veränderlichkeit im Laufe der ganzen Beobachtungsreihe zu erkennen, kann man nun zuerst den ersten Satz mit dem zweiten auf folgende Art vergleichen. Die Dauer von 147 Schwingungen findet sich aus
\begin{center}
\begin{tabular}{rcl}
$0\text{—}147$ & $.\ .\ .\ .\ .$ & $1^{\mathrm{h}}43'22''70$\\
$1\text{—}148$ & & $\phantom{1^{\mathrm{h}}43'}22.75$\\
$2\text{—}149$ & & $\phantom{1^{\mathrm{h}}43'}22.70$\\
$3\text{—}150$ & & $\phantom{1^{\mathrm{h}}43'}22.80$\\
$4\text{—}151$ & & $\phantom{1^{\mathrm{h}}43'}22.75$
\end{tabular}
\end{center}
im Mittel $1^{\mathrm{h}}43'22''74$ oder die Dauer Einer $=42''19551$. Auf gleiche Weise erhält man die Schwingungsdauer zwischen dem zweiten und dritten Satze $=42''17654$, zwischen dem dritten und vierten $=42''17879$, und zwischen dem ersten und vierten oder das Mittel aus der ganzen Reihe $=42''18344$.

Diese Rechnung kann auch in einer etwas abgeänderten Form geführt werden, die zugleich den Vortheil einer klaren Übersicht des regelmässigen Ganges sämmtlicher einzelnen Beobachtungen gewährt. Man fängt damit an, die einzelnen gefundenen Elongationszeiten mit einem genäherten Werthe der Schwingungsdauer auf einerlei Epoche zu reduciren, indem man von jeder den Betrag aller seit dieser Epoche verflossenen Schwingungszeiten, mit Hülfe dieses genäherten Werthes zurückrechnet. Man subtrahirt also von jeder Zahl der zweiten
\SourceLeaf{385}Columne das Product dieses angenommenen Werthes in die entsprechende Zahl der ersten Columne. Hätten die Beobachtungen eine absolute Genauigkeit, und wäre die Schwingungsdauer genau constant, und dem angenommenen Werthe genau gleich, so müssten sämmtliche so reducirte Zahlen genau gleich ausfallen. Aus dem Zunehmen der Zahlen von einem Satze zu dem folgenden hingegen erkennt man, dass die zum Grunde gelegte Schwingungsdauer für diesen Zeitraum zu klein war, und umgekehrt, während das unregelmässige Hinundherspringen der zu einem und demselben Satze gehörenden Resultate einen Maassstab für die Genauigkeit der Beobachtungen selbst darbietet.

In unserm Beispiele folgt aus der Vergleichung der ersten Elongationszeit mit der letzten die Schwingungsdauer $=42''18389$, anstatt welcher der genäherte Werth $42''18$ zur Berechnung der Zahlen der dritten Columne zum Grunde gelegt ist. Man sieht so mit Einem Blick, dass diese Schwingungsdauer für die Zeit vom ersten zum zweiten Satze etwas zu klein, hingegen von dem zweiten zum dritten. und eben so vom dritten zum vierten um ein geringes zu gross ist. Um genaue Resultate zu erhalten, nimmt man aus den zu jedem Satze gehörenden Zahlen der dritten Columne das Mittel; diese Mittel
\begin{center}
\begin{tabular}{r}
$21^{\mathrm{h}}55'47''69$\\
$49.97$\\
$49.52$\\
$49.35$
\end{tabular}
\end{center}
können als schärfere Werthe der zu den Ordnungszahlen 2, 149, 279, 420 gehörenden reducirten Zeiten angesehen werden. Man hat also vom ersten Satze zum zweiten ein Voreilen der Beobachtungen von $2''28$ vor dem vorausgesetzten Gange während 147 Schwingungen, was auf Eine Schwingung $0''01551$ beträgt, so dass der corrigirte Werth $42''19551$ wird, genau mit dem oben gefundenen übereinstimmend. Denselben Erfolg ergibt die Vergleichung der folgenden Sätze.

Für die Güte der Beobachtungen selbst gibt der blosse Anblick der zu einerlei Satz gehörenden Zahlen Zeugniss; indess mögen hier die Vorschriften Platz finden, wonach man in geeigneten Fällen den Maassstab für die Genauigkeit bestimmter ausmitteln kann. Bezeichnet man den mittlern bei einem Antritt zu befürchtenden Fehler mit $\varepsilon$, die Anzahl der zu einem Satze gehörenden Resultate mit $s$, und die Summe der Quadrate der Differenzen dieser einzelnen Resultate
\SourceLeaf{386}von ihrem Mittel mit $q$, so kann man näherungsweise annehmen
\[
\frac{(s-1)^2\varepsilon\varepsilon}{2s}=q
\]
oder wenn mehrere Sätze vorhanden sind,
\[
\varepsilon\varepsilon\sum\frac{(s-1)^2}{2s}=\sum q
\]
also
\[
\varepsilon=\BareRoot{\frac{\sum q}{\sum\frac{(s-1)^2}{2s}}}
\]

In unserm Beispiele sind bei dem ersten Satze die Differenzen vom Mittel $0''04$, $0''07$, $0''05$, $0''02$, $0''04$, also $q=110$, wenn man das Hunderttheil der Secunde als Einheit betrachtet; ferner $s=5$, also $\frac{(s-1)^2}{2s}=\frac85$. Für die drei folgenden Sätze ist, bei gleichem Werthe von $s$, $q=190;\ 110;\ 50$. Wir haben also
\[
\tfrac{32}{5}\varepsilon\varepsilon=460
\]
oder
\[
\varepsilon=8.5,\qquad\text{d. i. }\varepsilon=0''085.
\]

Indessen muss bemerkt werden, dass die Gültigkeit dieser Vorschrift von mehrern Bedingungen abhängig ist, die unserm Beispiele nicht hinlänglich eigen sind: erstlich nemlich, dass der vorausgesetzte genäherte Werth der Schwingungsdauer, womit die reducirten Zahlen berechnet sind, ohne merklichen Fehler als der wahre während jedes Satzes betrachtet werden dürfe; zweitens, dass die verschiedenen Sätze, die man vereinigt, unter nahe gleichen Umständen (so weit sie die Genauigkeit des Beobachtens afficiren können) beobachtet seien. Beides trifft in unserm Beispiel nicht zu, und man hat daher obige Rechnung nur wie eine Erläuterung der Formel zu betrachten. Will man genauere Bestimmungen haben, so ist es besser, zunächst zu diesem Zweck besondere Beobachtungen zu machen. Unter dem Vorbehalt, diesen Gegenstand in Zukunft ausführlicher zu behandeln, mag hier nur bemerkt werden, dass die Genauigkeit des Beobachtens — neben der Individualität des Apparats und des Beobachters — auch nach der bessern oder schlechtern Beleuchtung der Scale, der Schnelligkeit der Schwingungsbewegung, und der Beschaffenheit der Uhr ungleich ist. Eine gar zu schnelle Bewegung sowohl, als eine gar zu langsame ist der Genauigkeit des Be\SourceLeaf{387}obachtens weniger günstig, als eine mittlere Geschwindigkeit, und an einer Secundenuhr beobachtet man nicht so scharf, als an einem Chronometer, welches kleinere Zeittheile schlägt. Unter den günstigsten Umständen übertrifft die Genauigkeit dieser Beobachtungen sehr weit die der besten Beobachtungen an einem Mittagsfernrohr.

Die Schwingungsdauer ist bekanntlich, alles übrige gleich gesetzt, desto kleiner, je kleiner der Schwingungsbogen ist, und zwar so, dass während dieser sich dem Verschwinden unendlich nähert, jener einen Grenzwerth hat. Bezeichnen wir diesen Grenzwerth, oder, nach gewöhnlicher Sprachweise, die Zeit einer unendlich kleinen Schwingung, mit $T$, die einem Schwingungsbogen $G$ entsprechende Dauer hingegen mit $T'$, so hat man bekanntlich
\[
T'=T(1+\tfrac14\sin\tfrac14G^2+\tfrac14.\tfrac9{16}\sin\tfrac14G^4+\tfrac14.\tfrac9{16}.\tfrac{25}{36}\sin\tfrac14G^6+\text{etc.})
\]

Bei der Kleinheit der Bögen, auf welche man beim Gebrauch des Magnetometers beschränkt bleibt, kann man die Glieder der vierten und höhern Ordnung unbedenklich bei Seite, und deshalb auch $\frac{g}{8r}$ anstatt $\sin\tfrac14G$ setzen, wo $g$ das dem Bogen $G$ entsprechende Stück der Scale, und $r$ die horizontale Entfernung der Mitte der Scale vom Spiegel bedeutet. Wir haben also $T'=T(1+\frac{gg}{256rr})$, oder mit derselben Genauigkeit $T=T'(1-\frac{gg}{256rr})$. Für unser Beispiel ist in Scalentheilen oder Millimetern $r=4775.9$. Der Schwingungsbogen zwischen den Elongationen 0 und 1 ist $=1485.8$, und damit die Reduction der Schwingungsdauer auf eine unendlich kleine Schwingung $=-0''01595$. Eben so findet sich die Reduction der zweiten Schwingung $0''01590$, die der dritten $0''01584$, die der vierten $0''01577$, so dass im Mittel aus den vier ersten Schwingungen die auf unendlich kleine reducirte Dauer sich $=42''18414$ ergibt.

Die Anwendbarkeit dieses Verfahrens setzt aber die ununterbrochene Beobachtung der Elongation voraus. Die Reduction einer Reihe von Schwingungen, wovon nur Anfang und Ende beobachtet ist, auf unendlich kleine Bögen, mag man in dem Falle, wo der Schwingungsbogen in der Zwischenzeit nur eine mässige Abnahme erlitten hat, allenfalls so ausführen, dass man einen mittlern Werth der Grösse des Schwingungsbogens dabei zum Grunde legt. Allein die Reduction einer längern Reihe solcher Schwingungen, während welcher der Bogen sich stark vermindert hat, erfordert nothwendig eine wenigstens näherungsweise richtige
\SourceLeaf{388}Kenntniss des Gesetzes, nach welchem diese Verminderung geschieht. Ich gehe daher zu der Behandlung der beobachteten Elongationen über, deren nächste Resultate in der folgenden Tafel III. enthalten sind.
\begin{center}
\fontsize{9}{11}\selectfont
$\left.\begin{array}{r|r|r|r}
0&1009.725&1487.45&3.17244\\
1&1009.525&1484.55&3.17160\\
2&1009.425&1481.85&3.17081\\
3&1009.475&1478.85&3.16992\\
4&1009.075&1474.95&3.16878
\end{array}\right\}3.170710$\\[6pt]
$\left.\begin{array}{r|r|r|r}
147&1000.575&1003.25&3.00141\\
148&1000.375&1000.55&3.00024\\
149&1000.225&997.75&2.99902\\
150&1000.200&995.20&2.99791\\
151&1000.400&992.20&2.99660
\end{array}\right\}2.999036$\\[6pt]
$\left.\begin{array}{r|r|r|r}
277&994.050&694.90&2.84192\\
278&993.875&693.15&2.84083\\
279&993.700&691.40&2.83973\\
280&993.450&689.50&2.83853\\
281&993.350&687.30&2.83715
\end{array}\right\}3.839630$\\[6pt]
$\left.\begin{array}{r|r|r|r}
418&1003.725&455.65&2.65863\\
419&1003.575&454.85&2.65787\\
420&1003.125&453.45&2.65653\\
421&1002.950&451.50&2.65466\\
422&1002.925&449.85&2.65307
\end{array}\right\}2.656152$
\end{center}

Die erste Columne enthält die Ordnungszahl jeder Elongation; die zweite den entsprechenden Ruhestand der Nadel, nach der Formel $\tfrac14(a+2b+c)$, wenn $b$ die beobachtete Elongation, $a$ und $c$ die vorhergehende und folgende bedeuten (vergl. \textit{Resultate} für 1836, II); in der dritten Columne steht die doppelte Entfernung jeder Elongation von dem entsprechenden Ruhestande, oder der Bogen, welcher ohne die Ursachen, welche ihn zu vermindern streben, von da an beschrieben sein würde, also $\tfrac12(a-2b+c)$ oder $\tfrac12(2b-a-c)$, d. i. das Mittel des vorhergehenden und folgenden Schwingungsbogens; in der vierten Columne befindet sich der Logarithme dieser Zahl; endlich daneben der Mittelwerth aus den Zahlen der vierten Columne, die zu einem Satze gehören.

Alle Erfahrungen stimmen dahin überein, dass man, wenigstens während einer mässig grossen Zeit, die Zahlen der dritten Columne als in geometrischer,
\SourceLeaf{389}mithin ihre Logarithmen als in arithmetischer Progression abnehmend betrachten, wenigstens dies als die plausibelste Annäherung gelten lassen darf. Die kleinen Unregelmässigkeiten, welche sich bei der Vergleichung auf einander folgender Zahlen eines Satzes finden, hat man nur den unvermeidlichen kleinen Beobachtungsfehlern oder zufälligen kleinen Störungen zuzuschreiben, und man vermindert den nachtheiligen Einfluss davon, so viel thunlich, wenn man die Mittelzahlen in der fünften Columne als den entsprechenden mittlern Ordnungszahlen angehörig betrachtet, und daraus dann den Gang während der ganzen Beobachtungsreihe ableitet.

Wir haben demnach, als Logarithmen der Amplituden für die Elongationen
\begin{center}
\begin{tabular}{rcl}
$2$ & $.\ .\ .\ .\ .\ .\ .\ .\ .\ .$ & $3.170710$\\
$149$ & $.\ .\ .\ .\ .\ .\ .\ .\ .\ .$ & $2.999036$\\
$279$ & $.\ .\ .\ .\ .\ .\ .\ .\ .\ .$ & $2.839630$\\
$420$ & $.\ .\ .\ .\ .\ .\ .\ .\ .\ .$ & $2.656152$
\end{tabular}
\end{center}

Der Logarithme hat also vom ersten zum zweiten Satze während 147 Schwingungen die Abnahme $0.171674$ erlitten, was nach gleichförmiger Vertheilung auf Eine Schwingung $0.00116785$ beträgt: ich nenne diesen Quotienten das \textit{logarithmische Decrement}. Von dem zweiten zum dritten Satze findet sich dasselbe $=0.00122620$, vom dritten zum vierten $=0.00130126$. Man sieht, dass an einer gleichförmigen Abnahme hier wenigstens nicht viel fehlt: ich werde aber unten auf die Veränderlichkeit des logarithmischen Decrements zurückkommen.

Unter der Voraussetzung nun, dass die Amplituden während einer Reihe von Schwingungen in geometrischer Progression abgenommen haben, lassen sich diese auf unendlich kleine leicht reduciren. Ist $g$ die Grösse der ersten Schwingung in Scalentheilen, $g^0$ die Grösse der letzten, $\theta$ der Exponent der geometrischen Progression, also $g^0=g\theta^\mu$, wenn $\mu$ die Anzahl der Schwingungen bedeutet, so wird die Reduction der ersten Schwingungszeit auf die unendlich kleine Schwingung
\[
-\frac{Tgg}{256rr},
\]
die der zweiten
\[
=-\frac{Tgg\theta\theta}{256rr}
\]
u. s. w. also die Summe aller
\[
=-\frac{T(gg-g^0g^0\theta\theta)}{256rr(1-\theta\theta)}
\]

\SourceLeaf{390}Bezeichnen wir die zu der Anfangs- und End-Elongation gehörenden Amplituden, nach derselben Art berechnet wie in Tafel III, mit $h$ und $h^0$, so ist $h=\tfrac12(\frac{g}{\theta}+g)$, $h^0=\tfrac12(g^0+g^0\theta)$, mithin obige Summe
\[
=-\frac{T(hh-h^0h^0)}{64rr}.\frac{\theta\theta}{(1+\theta)^2(1-\theta\theta)}
\]

Da in allen hier in Rede stehenden Fällen $\theta$ ein von der Einheit wenig verschiedener, also das mit $\lambda$ zu bezeichnende logarithmische Decrement $=\log\frac1\theta$ ein kleiner Bruch ist, so kann man anstatt des zweiten Factors in jenem Ausdruck, für welchen sich, $m$ den Modulus des Logarithmensystems bedeutend, folgende Reihe ergibt:
\[
\frac{\theta\theta}{(1+\theta)^2(1-\theta\theta)}=\frac{m}{8\lambda}-\tfrac5{96}.\frac{\lambda}{m}+\tfrac{37}{2880}(\frac{\lambda}{m})^3\ \text{etc.}
\]
mit hinlänglicher Schärfe bloss das erste Glied $\frac{m}{8\lambda}$ setzen. Die Reduction der ganzen Dauer der $\mu$ Schwingungen wird also:
\[
=-\frac{Tm(hh-h^0h^0)}{512rr\lambda}
\]
oder die Reduction des Durchschnittwerths für Eine Schwingung
\[
=-\frac{Tm(hh-h^0h^0)}{512rr\lambda\mu}
\]

Es ist bei diesen Formeln aus den oben angeführten Gründen gleichgültig, ob man darin für $T$ den berichtigten oder den unberichtigten Werth gebraucht. In unserm Beispiele findet sich der Werth der Reduction
\begin{center}
\fontsize{9}{11}\selectfont
\begin{tabular}{r|r|r|r}
& \multicolumn{1}{c|}{für die ganze} & \multicolumn{1}{c|}{für Eine} & \multicolumn{1}{c}{Reducirter}\\
& \multicolumn{1}{c|}{Zeit} & \multicolumn{1}{c|}{Schwingung} & \multicolumn{1}{c}{Werth}\\[5pt]
$2\ .\ .\ .\ .\ 149$ & $-1''6111$ & $-0''01096$ & $42''18455$\\
$149\ .\ .\ .\ .\ 279$ & $-0.6624$ & $-0.00510$ & $42.17144$\\
$279\ .\ .\ .\ .\ 420$ & $-0.3285$ & $-0.00233$ & $42.17646$
\end{tabular}
\end{center}

Die Abnahme des Schwingungsbogens ist, auch ausser ihrem Zusammenhange mit der genauern Berechnung der Schwingungsdauer, noch in mehrern andern Beziehungen von Interesse. Man hat dabei zunächst die äussern Umstände zu unterscheiden, unter welchen die Nadel ihre Schwingungen macht.

\SourceLeaf{391}Wenn der Apparat zweckmässig eingerichtet und in vollkommen gutem Zustande ist\HistoricalNote{391}, und in seinen nächsten Umgebungen sich Nichts befindet, was eine beträchtliche Dämpfung der Schwingungsbewegung bewirken muss, so ist die Abnahme der Schwingungsbögen immer sehr langsam, und in so fern regelmässig, als sie wenigstens während einer mässigen Zeit in geometrischer Progression erfolgt, mithin das logarithmische Decrement nahe constant ist. In unserm Beispiele ändert sich dieses logarithmische Decrement während fast sechs Stunden nur von $0.00117$ bis $0.00130$, oder die Abnahme des Bogens von einer Schwingung zur andern schwankte von $\tfrac1{371}$ bis $\tfrac1{334}$. Allein die Erfahrung zeigt, dass sehr häufig sehr verschiedene Werthe des logarithmischen Decrements vorkommen: es steigt wohl an demselben Apparate bis gegen $0.00300$, und sinkt zu andern Zeiten auf $0.00030$ und selbst, in seltenen Fällen, noch tiefer herab. Immer aber geschehen, nach unsern Erfahrungen, die Veränderungen nur allmählig. Man wird also diese Abnahme nicht wohl allein dem Widerstande der Luft zuschreiben dürfen: aber die eigentliche Ursache, welche diese Verschiedenheiten bedingt, hat sich bisher unserm vielfach wiederholten Versuchen entzogen, und wir wünschen daher sehr, dass auch Beobachter an andern Orten ihre besondere Aufmerksamkeit auf dieses zur Zeit noch räthselhafte Phänomen richten mögen. Ein Umstand ist bei diesen Versuchen so oft bemerkt, dass wir ihn kaum noch für zufällig halten können, wenn auch ein Causalnexus noch ganz unerklärlich bleibt, nemlich dass die \textit{sehr} kleinen Werthe des logarithmischen Decrements immer nur bei bedecktem, die sehr grossen hingegen gewöhnlich nur bei heiterm Himmel eintreten, wobei zum Überfluss noch bemerkt werden mag, dass der Apparat nicht an einem Seidenfaden, sondern an einem Metalldraht aufgehängt ist, und dass diese Versuche immer nur an windstillen Tagen angestellt, folglich in diesem beiden Beziehungen sowohl hygrometrischer Einfluss als Luftzug ganz ausser Frage sind.

Ist hingegen die Nadel von einem Multiplicator umgeben, der einen Theil einer geschlossenen Kette ausmacht, so tritt eine neue Ursache der Abnahme der Schwingungsbögen hinzu. Die Bewegung der Nadel inducirt nemlich in dieser Kette einen galvanischen Strom, dessen Intensität am stärksten ist, wenn die
